\chapter{The reciprocal zeta function}\label{chap:rzeta}

The reciprocal zeta function $1/\zeta(\frac12+\sqrt s)$ is, up to the factor $1-4s$ of the pole and a constant, the reciprocal Xi function times the gamma factor, and unlike $\xi$ it has a Dirichlet series and an Euler product, from which we read its Thorin type measures.

Sections~\ref{ch:rz1} and \ref{ch:rz2} compute the Thorin measure of the gamma factor, the Thorin measures of the normalised reciprocal zeta function and of the reciprocal Xi function under RH, and, unconditionally, the formal Thorin measures of these two functions in the Vere-Jones form of Section~\ref{rz:sec:lvj}. After the stable tilt $e^{-c\sqrt s}$, with $c=\zeta'(\frac12)/\zeta(\frac12)$, the reciprocal of the gamma factor is a generalised gamma convolution, and multiplying the normalised reciprocal zeta function by it gives the floor $e^{-c\sqrt s}\,\xi(\frac12)/\xi(\frac12+\sqrt s)$ of the reciprocal Xi function (Section~\ref{rz:sec:floor}), in which the zeros are the only possible source of negativity.

The infinite divisibility of the floor, as proved here, rests on the positivity of the canonical function of the normalised reciprocal zeta function, which is Theorem~\ref{rz:thm:id} of Chapter~\ref{chap:rzid}. It also follows from the infinite divisibility of the clock (Theorem~\ref{thm:ch12:idsd}); what the floor adds is a splitting of its canonical function into an arithmetic part and a completely monotone archimedean part. The floor is a generalised gamma convolution exactly when RH holds, and the derivatives of its canonical function alternate in sign up to order $10^{10}$ by the methods of this chapter and up to order $1.16\times10^{25}$ by the estimates of Section~\ref{ch:heat}.
